% Part: propositional-logic

\documentclass[../../include/open-logic-part]{subfiles}

\begin{document}

\olpart{pl}{Lógica Proposicional}

\begin{editorial}
  Esta parte contém material sobre lógica proposicional clássica. O
  primeiro capítulo é relativamente elementar e apenas enumera definições
  e resultados; muitas demonstrações não são apresentadas, mas deixadas como
  exercícios. O material sobre sistemas de demonstração e o teorema da
  completude é incorporado da parte sobre lógica de primeira ordem, com a
  etiqueta ``FOL'' definida como falsa. Isso exclui tudo o que diz respeito a
  predicados, termos e quantificadores, e substitui menções a
  !!{structure}s~$\Struct{M}$ por menções a
  !!{valuation}s~$\pAssign{v}$.

  Está previsto ampliar esta parte para incluir
  mais detalhes e acrescentar outros tópicos e resultados, como a
  completude verofuncional.
\end{editorial}

\olimport[syntax-and-semantics]{syntax-and-semantics}

\tagfalse{FOL}

\olimport[../first-order-logic/proof-systems]{proof-systems}

\iftag{prfSC}{%
  \olimport[../first-order-logic/sequent-calculus]{sequent-calculus}
}{}

\iftag{prfND}{%
  \olimport[../first-order-logic/natural-deduction]{natural-deduction}
}{}

\iftag{prfTab}{%
  \olimport[../first-order-logic/tableaux]{tableaux}
}{}

\iftag{prfAX}{%
  \olimport[../first-order-logic/axiomatic-deduction]{axiomatic-deduction}
}{}

\olimport[../first-order-logic/completeness]{completeness}

\tagtrue{FOL}

\OLEndPartHook

\end{document}
